\subsubsection{Advanced Matrix Analysis}

This subsection collects matrix identities that organise larger linear systems by blocks, tensor products, vectorisation, and rank bookkeeping.

\begin{definition} \label{def:kronecker-product} \hlt{Kronecker Product}\\
Let $A=[a_{ij}]\in\operatorname{Mat}(m,n,\mathbb{F})$ and $B\in\operatorname{Mat}(p,q,\mathbb{F})$. The \hlt{Kronecker product} of $A$ and $B$ is the block matrix
\begin{align}
A\otimes B
\defeq
\begin{bmatrix}
a_{11}B & a_{12}B & \cdots & a_{1n}B\\
a_{21}B & a_{22}B & \cdots & a_{2n}B\\
\vdots & \vdots & \ddots & \vdots\\
a_{m1}B & a_{m2}B & \cdots & a_{mn}B
\end{bmatrix}
\in\operatorname{Mat}(mp,nq,\mathbb{F}).
\nonumber
\end{align}
The Kronecker product is also called the tensor product of matrices.
\end{definition}

\begin{proposition} \label{prop:basic-kronecker-identities} \hlt{Basic Kronecker Identities}\\
Whenever the expressions are defined,
\begin{enumerate}[label=(\roman*)]
\setlength{\itemsep}{0pt}
\item $c(A\otimes B)=(cA)\otimes B=A\otimes(cB)$;
\item $(A+B)\otimes C=A\otimes C+B\otimes C$ and $A\otimes(B+C)=A\otimes B+A\otimes C$;
\item $(A\otimes B)\otimes C=A\otimes(B\otimes C)$;
\item $(A\otimes B)^t=A^t\otimes B^t$;
\item $I_m\otimes I_n=I_{mn}$.
\end{enumerate}
If $\mathbb{F}=\mathbb{C}$, then also $\overline{A\otimes B}=\overline{A}\otimes\overline{B}$ and $(A\otimes B)^*=A^*\otimes B^*$.
\begin{proof}
\begin{enumerate}[label=(\roman*)]
\setlength{\itemsep}{0pt}
\item The $(i,j)$ block of $c(A\otimes B)$ is $c(a_{ij}B)=(ca_{ij})B=a_{ij}(cB)$, which is $(i,j)$ block of $(cA)\otimes B$ and $A\otimes(cB)$.
\item The $(i,j)$ block of $(A+B)\otimes C$ is $(a_{ij}+b_{ij})C=a_{ij}C+b_{ij}C$, so $(A+B)\otimes C=A\otimes C+B\otimes C$.\\
The proof of $A\otimes(B+C)=A\otimes B+A\otimes C$ is identical inside each block.
\item If $B=[b_{rs}]$, then both $(A\otimes B)\otimes C$ and $A\otimes(B\otimes C)$ have blocks $a_{ij}b_{rs}C$ in the same relative positions.
\item The $(j,i)$ block of $(A\otimes B)^t$ is $(a_{ij}B)^t=a_{ij}B^t$, which is the $(j,i)$ block of $A^t\otimes B^t$.
\item The block matrix $I_m\otimes I_n$ has $I_n$ on its diagonal blocks and zero elsewhere, hence equals $I_{mn}$.
\end{enumerate}
If $\mathbb{F}=\mathbb{C}$, conjugation is checked blockwise, adjoint identity follows from conjugation together with part (iv).\\
\hlt{Structural Commentary}: Kronecker products replace scalar entries by whole scaled blocks. Standard algebraic rules survive because each block is still controlled by scalar distributivity and transposition.\\
\hlt{Key Insight}: Check identities block-by-block; each block reduces to an ordinary scalar-matrix identity.
\end{proof}
\end{proposition}

\begin{theorem} \label{thm:kronecker-mixed-product-property} \hlt{Mixed Product Property}\\
Let $A\in\operatorname{Mat}(m,n,\mathbb{F})$, $B\in\operatorname{Mat}(p,q,\mathbb{F})$, $C\in\operatorname{Mat}(n,r,\mathbb{F})$, $D\in\operatorname{Mat}(q,s,\mathbb{F})$.  $(A\otimes B)(C\otimes D)=(AC)\otimes(BD)$.
\begin{proof}
Write $A=[a_{ik}]$ and $C=[c_{kj}]$.\\
The $(i,j)$ block of $(A\otimes B)(C\otimes D)$ is $\sum\limits_{k=1}^{n}(a_{ik}B)(c_{kj}D) =\left(\sum\limits_{k=1}^{n}a_{ik}c_{kj}\right)BD =(AC)_{ij}BD$, which is exactly the $(i,j)$ block of $(AC)\otimes(BD)$.\\
\hlt{Structural Commentary}: The outer indices multiply as $AC$, while the inner block action multiplies as $BD$.\\
\hlt{Key Insight}: Kronecker multiplication separates into ordinary multiplication in each factor.
\end{proof}
\end{theorem}

\begin{corollary} \label{cor:inverse-of-kronecker-product} \hlt{Inverse of a Kronecker Product}\\
If $A\in\operatorname{Mat}(m,m,\mathbb{F})$ and $B\in\operatorname{Mat}(n,n,\mathbb{F})$ are invertible, then $A\otimes B$ is invertible and $(A\otimes B)^{-1}=A^{-1}\otimes B^{-1}$.
\begin{proof}
By Theorem \ref{thm:kronecker-mixed-product-property}, $(A\otimes B)(A^{-1}\otimes B^{-1})=I_m\otimes I_n=I_{mn}$ and the reverse product is identical.\\
\hlt{Structural Commentary}: Inverting a Kronecker product inverts each independent factor.\\
\hlt{Key Insight}: The mixed product property turns the inverse check into two ordinary inverse checks.
\end{proof}
\end{corollary}

\begin{definition} \label{def:vectorization-map} \hlt{Vectorisation}\\
Let $X=[\mathbf{x}_1\ \mathbf{x}_2\ \cdots\ \mathbf{x}_n]\in\operatorname{Mat}(m,n,\mathbb{F})$. The \hlt{vectorisation} of $X$ is
\begin{align}
\operatorname{vec}(X)\defeq
\begin{bmatrix}
\mathbf{x}_1\\
\mathbf{x}_2\\
\vdots\\
\mathbf{x}_n
\end{bmatrix}
\in\mathbb{F}^{mn}.
\nonumber
\end{align}
Thus $\operatorname{vec}$ stacks the columns of $X$ vertically.
\end{definition}

\begin{theorem} \label{thm:vectorization-of-triple-product} \hlt{Vectorisation of a Triple Product}\\
Let $A\in\operatorname{Mat}(m,n,\mathbb{F})$, $X\in\operatorname{Mat}(n,p,\mathbb{F})$, and $B\in\operatorname{Mat}(p,q,\mathbb{F})$. Then $\operatorname{vec}(AXB)=(B^t\otimes A)\operatorname{vec}(X)$.
\begin{proof}
Write $B=[\mathbf{b}_1\ \cdots\ \mathbf{b}_q]$ and $X=[\mathbf{x}_1\ \cdots\ \mathbf{x}_p]$.\\
The $k$th column of $AXB$ is $AX\mathbf{b}_k=A\sum\limits_{i=1}^{p}b_{ik}\mathbf{x}_i =\sum\limits_{i=1}^{p}b_{ik}A\mathbf{x}_i =(\mathbf{b}_k^t\otimes A)\operatorname{vec}(X)$.\\
Stacking these $q$ columns gives $(B^t\otimes A)\operatorname{vec}(X)$.\\
\hlt{Structural Commentary}: Vectorisation rewrites two-sided multiplication as one matrix action.\\
\hlt{Key Insight}: Left multiplication contributes $A$; right multiplication contributes $B^t$.
\end{proof}
\end{theorem}

\begin{theorem} \label{thm:sylvester-linear-matrix-equation} \hlt{Sylvester Linear Matrix Equation}\\
Let $A\in\operatorname{Mat}(m,m,\mathbb{C})$ and $B\in\operatorname{Mat}(n,n,\mathbb{C})$. If $\operatorname{spec}(A)\cap\operatorname{spec}(B)=\emptyset$, then for every $C\in\operatorname{Mat}(m,n,\mathbb{C})$ the equation $AX-XB=C$ has a unique solution $X\in\operatorname{Mat}(m,n,\mathbb{C})$.
\begin{proof}
Let $L(X)=AX-XB$. Since $L$ has equal finite-dimensional domain and codomain, it suffices to show $\ker L=\{0\}$. If $AX=XB$, then $p(A)X=Xp(B)$ for every polynomial $p$. \\
With $p=p_B$, Cayley--Hamilton gives $p_B(A)X=Xp_B(B)=0$. For every $\lambda\in\operatorname{spec}(A)$, disjointness gives $p_B(\lambda)\neq0$; hence $p_B(A)$ is invertible by Corollary \ref{cor:polynomial-in-matrix-invertible-spectrum-test}. Thus $X=0$, so $L$ is bijective.\\
\hlt{Structural Commentary}: Polynomial calculus pass through homogeneous solutions; disjoint spectra invert $p_B(A)$.\\
\hlt{Key Insight}: Spectral separation turns an intertwining relation into the zero relation.
\end{proof}
\end{theorem}

\begin{figure}[H]
  \centering
  \begin{tikzpicture}[
  >=Stealth,
  every node/.style={font=\small},
  box/.style={draw=gray!55, fill=gray!5, rounded corners=2pt, align=center, inner xsep=5pt, inner ysep=4.5pt},
  good/.style={draw=blue!60!black, fill=blue!7, rounded corners=2pt, align=center, inner xsep=5pt, inner ysep=4.5pt},
  arr/.style={->, line width=0.75pt, blue!70!black}
]
  \node[box] (gap) at (0,0) {$\operatorname{spec}(A)\cap\operatorname{spec}(B)=\emptyset$};
  \node[good] (poly) at (3.61,0) {$p_B(A)$ invertible};
  \node[box] (map) at (6.78,0) {$L^{-1}$ exists};

  \draw[arr] (gap.east) -- (poly.west);
  \draw[arr] (poly.east) -- (map.west);
\end{tikzpicture}

  \caption{Disjoint spectra make $p_B(A)$ invertible, so the Sylvester map $L(X)=AX-XB$ is invertible.}
  \label{fig:sylvester-equation-spectral-gap}
\end{figure}

\begin{corollary} \label{cor:block-upper-triangular-remove-disjoint-coupling} \hlt{Removing a Spectrally Separated Off-Diagonal Block}\\
Let $M=\begin{bmatrix}A&C\\0&B\end{bmatrix}$ and $\operatorname{spec}(A)\cap\operatorname{spec}(B)=\emptyset$ Then $M$ is similar to $A\oplus B$.
\begin{proof}
By Theorem \ref{thm:sylvester-linear-matrix-equation}, choose $X$ such that $AX-XB=-C$. Then
\begin{align}
\begin{bmatrix}I&-X\\0&I\end{bmatrix}
\begin{bmatrix}A&C\\0&B\end{bmatrix}
\begin{bmatrix}I&X\\0&I\end{bmatrix}
=
\begin{bmatrix}A&AX+C-XB\\0&B\end{bmatrix}
=
\begin{bmatrix}A&0\\0&B\end{bmatrix}.
\nonumber
\end{align}
\hlt{Structural Commentary}: A block shear cancels the off-diagonal block via Sylvester's equation.\\
\hlt{Key Insight}: Spectrally separated diagonal blocks cannot carry essential upper-right coupling.
\end{proof}
\end{corollary}

\begin{theorem} \label{thm:schur-block-form-by-eigenvalue-groups} \hlt{Schur Form Grouped by Eigenvalues}\\
Let $A\in\operatorname{Mat}(n,n,\mathbb{C})$ have distinct eigenvalues $\lambda_1,\ldots,\lambda_d$ with algebraic multiplicities $n_1,\ldots,n_d$.\\
Then $A$ is unitarily similar to a block upper-triangular matrix
\begin{align}
T=
\begin{bmatrix}
T_{11}&T_{12}&\cdots&T_{1d}\\
0&T_{22}&\cdots&T_{2d}\\
\vdots&\vdots&\ddots&\vdots\\
0&0&\cdots&T_{dd}
\end{bmatrix},
\qquad
T_{ii}\in\operatorname{Mat}(n_i,n_i,\mathbb{C}),
\nonumber
\end{align}
where each $T_{ii}$ is upper triangular with all diagonal entries equal to $\lambda_i$. Moreover, $A$ is similar to $T_{11}\oplus\cdots\oplus T_{dd}$.
\begin{proof}
Order the Schur diagonal as $\lambda_1,\ldots,\lambda_1,\lambda_2,\ldots,\lambda_2,\ldots,\lambda_d,\ldots,\lambda_d$.\\
Theorem \ref{thm:matrix-schur-triangularization-ordered} gives the displayed block Schur form. Since $\operatorname{spec}(T_{11})$ is disjoint from the trailing block spectrum, Corollary \ref{cor:block-upper-triangular-remove-disjoint-coupling} splits off $T_{11}$. Iterating splits off every diagonal block.\\
\hlt{Structural Commentary}: Grouped Schur form isolates spectral blocks; Sylvester shears split disjoint groups.\\
\hlt{Key Insight}: Cross-group upper blocks are coordinate artefacts.
\end{proof}
\end{theorem}

\begin{theorem} \label{thm:minimal-polynomial-direct-sum-disjoint-spectra} \hlt{Minimal Polynomial of a Spectrally Disjoint Direct Sum}\\
Let $A=A_1\oplus\cdots\oplus A_d$, where the spectra of $A_1,\ldots,A_d$ are pairwise disjoint. Then $\mu_A=\mu_{A_1}\mu_{A_2}\cdots\mu_{A_d}$.
\begin{proof}
Let $f=\prod\limits_i\mu_{A_i}$. Since $f(A_i)=0$ for all $i$, $f(A)=0$ and $\mu_A\mid f$. Conversely, $\mu_A(A_i)=0$, so $\mu_{A_i}\mid\mu_A$ for each $i$. Disjoint spectra make the $\mu_{A_i}$ pairwise relatively prime, hence $f\mid\mu_A$. Thus $\mu_A=f$.\\
\hlt{Structural Commentary}: Direct sums are annihilated blockwise; disjoint spectra make the block factors coprime.\\
\hlt{Key Insight}: Each spectral component contributes its own minimal-polynomial factor.
\end{proof}
\end{theorem}

\begin{corollary} \label{cor:commuting-with-disjoint-blocks-has-zero-off-diagonal-blocks} \hlt{Commuting with Disjoint Blocks}\\
Let $A=A_{11}\oplus\cdots\oplus A_{dd}$, where the spectra of $A_{11},\ldots,A_{dd}$ are pairwise disjoint.\\
Partition $B=[B_{ij}]$ conformally with $A$. If $AB=BA$, then $B_{ij}=0$ for all $i\neq j$.
\begin{proof}
Equating $(i,j)$ blocks in $AB=BA$ gives $A_{ii}B_{ij}-B_{ij}A_{jj}=0$.\\
For $i\neq j$, the spectra of $A_{ii}$ and $A_{jj}$ are disjoint, so Theorem \ref{thm:sylvester-linear-matrix-equation} with $C=0$ gives $B_{ij}=0$.\\
\hlt{Structural Commentary}: A commuting matrix cannot mix direct-sum blocks with disjoint spectra.\\
\hlt{Key Insight}: The off-diagonal block solves a homogeneous Sylvester equation, whose only solution is zero.
\end{proof}
\end{corollary}

\begin{definition} \label{def:nonderogatory-matrix} \hlt{Nonderogatory Matrix}\\
A matrix $A\in\operatorname{Mat}(n,n,\mathbb{C})$ is \hlt{nonderogatory} if its Jordan form has exactly one Jordan block for each distinct eigenvalue. Equivalently, $\dim E_{\lambda}(A)=1$ for every $\lambda\in\operatorname{spec}(A)$.
\end{definition}

\begin{lemma} \label{lem:single-jordan-block-commutant-polynomials} \hlt{Commutant of One Jordan Block}\\
If $J=J_m(\lambda)$ and $XJ=JX$, then $X=p(J)$ for some polynomial $p$ with $\deg p<m$.
\begin{proof}
It suffices to treat $N=J_m(0)$, since $J_m(\lambda)=\lambda I_m+N$. The vector $\mathbf{e}_m$ is cyclic for $N$:
\begin{align}
(\mathbf{e}_m,N\mathbf{e}_m,\ldots,N^{m-1}\mathbf{e}_m)
\nonumber
\end{align}
is the standard basis in reverse order. Choose the unique polynomial $p$ with $\deg p<m$ such that $p(N)\mathbf{e}_m=X\mathbf{e}_m$. For every $r=0,\ldots,m-1$,
\begin{align}
p(N)N^r\mathbf{e}_m=N^rp(N)\mathbf{e}_m=N^rX\mathbf{e}_m=XN^r\mathbf{e}_m.
\nonumber
\end{align}
Thus $p(N)$ and $X$ agree on a basis, so $X=p(N)=q(J)$ for $q(z)=p(z-\lambda)$.\\
\hlt{Structural Commentary}: A single Jordan block is generated by one cyclic vector.\\
\hlt{Key Insight}: A commuting matrix is determined by where it sends the cyclic vector, hence by a polynomial in the block.
\end{proof}
\end{lemma}

\begin{theorem} \label{thm:nonderogatory-commutant-polynomials} \hlt{Nonderogatory Commutants Are Polynomials}\\
If $A\in\operatorname{Mat}(n,n,\mathbb{C})$ is nonderogatory, then
\begin{align}
\{A\}'=\{p(A):p\in\mathcal{P}(\mathbb{C})\}.
\nonumber
\end{align}
\begin{proof}
The inclusion $\{p(A):p\in\mathcal{P}(\mathbb{C})\}\subseteq\{A\}'$ is immediate.\\
For the reverse inclusion, replace $A$ by its Jordan form
\begin{align}
J=J_{m_1}(\lambda_1)\oplus\cdots\oplus J_{m_d}(\lambda_d),
\nonumber
\end{align}
where the $\lambda_i$ are distinct. If $XJ=JX$ and $X=[X_{ij}]$ is partitioned conformally with $J$, Corollary \ref{cor:commuting-with-disjoint-blocks-has-zero-off-diagonal-blocks} gives $X_{ij}=0$ for $i\neq j$. Each diagonal block $X_{ii}$ commutes with $J_{m_i}(\lambda_i)$, so Lemma \ref{lem:single-jordan-block-commutant-polynomials} gives $X_{ii}=p_i(J_{m_i}(\lambda_i))$ for some $\deg p_i<m_i$.\\
Because the polynomials $(z-\lambda_i)^{m_i}$ are pairwise coprime, Hermite interpolation gives a polynomial $p$ such that $p\equiv p_i\pmod{(z-\lambda_i)^{m_i}}$ for every $i$. Then $X=p(J)$. Conjugating back to $A$ proves the claim.\\
\hlt{Structural Commentary}: Disjoint eigenvalue blocks force block diagonality; one Jordan block per eigenvalue makes each diagonal commutant polynomial.\\
\hlt{Key Insight}: Hermite interpolation glues the separate block polynomials into one polynomial in $A$.
\end{proof}
\end{theorem}

\begin{figure}[H]
  \centering
  \begin{tikzpicture}[
  font=\small,
  block/.style={draw=blue!60!black, fill=blue!6, rounded corners=2pt, align=center, inner sep=5pt, minimum width=2.00cm},
  poly/.style={draw=red!60!black, fill=red!6, rounded corners=2pt, align=center, inner sep=5pt, minimum width=2.33cm},
  arr/.style={-{Stealth[length=2mm]}, thick, gray!75}
]
  \node[block] (j1) at (0,0) {$J_{m_1}(\lambda_1)$};
  \node[block] (j2) at (2.55,0) {$J_{m_2}(\lambda_2)$};
  \node[block] (jd) at (5.10,0) {$\cdots$};
  \node[block] (jk) at (7.65,0) {$J_{m_d}(\lambda_d)$};

  \node[poly] (p1) at (0,-1.30) {$p_1(J_{m_1}(\lambda_1))$};
  \node[poly] (p2) at (2.55,-1.30) {$p_2(J_{m_2}(\lambda_2))$};
  \node[poly] (pd) at (5.10,-1.30) {$\cdots$};
  \node[poly] (pk) at (7.65,-1.30) {$p_d(J_{m_d}(\lambda_d))$};

  \node[poly, minimum width=3.52cm] (p) at (3.825,-2.52) {$p\equiv p_i\pmod{(z-\lambda_i)^{m_i}}$};

  \draw[arr] (j1) -- (p1);
  \draw[arr] (j2) -- (p2);
  \draw[arr] (jk) -- (pk);
  \draw[arr] (p1) -- (p);
  \draw[arr] (p2) -- (p);
  \draw[arr] (pk) -- (p);
\end{tikzpicture}

  \caption{For a nonderogatory matrix, commuting matrices are blockwise polynomials, then one interpolating polynomial glues the blocks.}
  \label{fig:nonderogatory-commutant-polynomial}
\end{figure}

\begin{theorem} \label{thm:simultaneous-schur-triangularization-commuting-family} \hlt{Simultaneous Schur Triangularisation}\\
Let $\emptyset\neq\mathcal{F}\subseteq\operatorname{Mat}(n,n,\mathbb{C})$ be a commuting family.\\
Then there is a unitary $U\in\operatorname{Mat}(n,n,\mathbb{C})$ such that $U^*AU$ is upper triangular for every $A\in\mathcal{F}$.\\
If every matrix in $\mathcal{F}$ is real and has only real eigenvalues, then $U$ can be chosen real orthogonal.
\begin{proof}
Induct on $n$. Choose a common unit eigenvector $\mathbf{x}$, extend it to an orthonormal basis, and write $V=[\mathbf{x}\ V_2]$. Then for every $A\in\mathcal{F}$,
\begin{align}
V^*AV=
\begin{bmatrix}
\lambda_A&*\\
0&A'
\end{bmatrix}.
\nonumber
\end{align}
Commutation gives $A'B'=B'A'$ for all $A,B\in\mathcal{F}$, so induction simultaneously triangularises the compressed family.\\
In the real case, choose a real common eigenvector: if all matrices are scalar, any real unit vector works; otherwise a nonzero proper real eigenspace of some nonscalar $A_0\in\mathcal{F}$ is invariant under all of $\mathcal{F}$, and induction inside it gives one. Then use real orthonormal extensions.\\
\hlt{Structural Commentary}: A common eigenline creates the first block; commutation descends to the compression.\\
\hlt{Key Insight}: Simultaneous triangularisation is Schur induction with a common eigenvector at every step.
\end{proof}
\end{theorem}

\begin{figure}[H]
  \centering
  \begin{tikzpicture}[
  >=Stealth,
  every node/.style={font=\small},
  box/.style={draw=gray!55, fill=gray!5, rounded corners=2pt, align=center, inner xsep=5pt, inner ysep=4pt},
  hi/.style={draw=blue!60!black, fill=blue!7, rounded corners=2pt, align=center, inner xsep=5pt, inner ysep=4pt},
  arr/.style={->, line width=0.75pt, blue!70!black}
]
  \node[hi] (common) at (0,0)
    {$\mathbb{C}\mathbf{x}$};
  \node[box] (block) at (3.18,0)
    {$V^*AV=\begin{bmatrix}\lambda_A&*\\0&A'\end{bmatrix}$};
  \node[box] (comm) at (6.60,0)
    {$A'B'=B'A'$};

  \draw[arr] (common.east) -- (block.west);
  \draw[arr] (block.east) -- (comm.west);
\end{tikzpicture}

  \caption{A shared eigenline gives a simultaneous first block; lower-right compressions still commute.}
  \label{fig:simultaneous-triangularization-commuting}
\end{figure}

\begin{corollary} \label{cor:commuting-sum-product-eigenvalues} \hlt{Eigenvalues of Sums and Products of Commuting Matrices}\\
Let $A,B\in\operatorname{Mat}(n,n,\mathbb{C})$ and $AB=BA$. Their eigenvalues can be ordered as $\lambda_1,\ldots,\lambda_n$ and $\mu_1,\ldots,\mu_n$ so that $A+B$ has eigenvalues $\lambda_i+\mu_i$ and $AB$ has eigenvalues $\lambda_i\mu_i$, counted with multiplicity.
\begin{proof}
Choose simultaneous Schur forms $T=U^*AU$ and $R=U^*BU$. Set $\lambda_i=t_{ii}$ and $\mu_i=r_{ii}$. Then $T+R$ and $TR$ are upper triangular with diagonals $\lambda_i+\mu_i$ and $\lambda_i\mu_i$, and are similar to $A+B$ and $AB$.\\
\hlt{Structural Commentary}: Simultaneous Schur form aligns both spectra on one diagonal.\\
\hlt{Key Insight}: In common triangular coordinates, diagonals add and multiply.
\end{proof}
\end{corollary}

\begin{lemma} \label{lem:adjoints-of-finitely-many-normal-matrices-polynomial} \hlt{Adjoints of Normal Matrices as One Polynomial}\\
Let $A_1,\ldots,A_k$ be normal complex matrices. Then $\exists$ a polynomial $p\in\mathcal{P}(\mathbb{C})$ s.t. $p(A_j)=A_j^* \ \forall j=1,\ldots,k$.
\begin{proof}
Let $\Omega=\bigcup\limits_{j=1}^{k}\operatorname{spec}(A_j)$. Since $\Omega$ is finite, choose a Lagrange interpolation polynomial $p$ satisfying $p(\lambda)=\overline{\lambda}$ for every $\lambda\in\Omega$. For each $A_j$, Corollary \ref{cor:normal-spectral-resolution-functional-calculus} gives $p(A_j)=A_j^*$.\\
\hlt{Structural Commentary}: One polynomial interpolates conjugation on all spectra.\\
\hlt{Key Insight}: For normal matrices, $A^*$ is $p(A)$ with $p(\lambda)=\overline{\lambda}$.
\end{proof}
\end{lemma}

\begin{theorem} \label{thm:fuglede-putnam-normal-intertwining} \hlt{Fuglede--Putnam Theorem}\\
Let $A\in\operatorname{Mat}(m,m,\mathbb{C})$ and $B\in\operatorname{Mat}(n,n,\mathbb{C})$ be normal, and let $X\in\operatorname{Mat}(m,n,\mathbb{C})$.\\
Then $AX=XB$ $\Longleftrightarrow$ $A^*X=XB^*$.
\begin{proof}
\begin{enumerate}[label=\roman*.]
\setlength{\itemsep}{0pt}
\item[$(\Rightarrow)$] Assume $AX=XB$. Then $A^rX=XB^r$ for every $r\geq0$, so $p(A)X=Xp(B)$ for every polynomial $p$. By Lemma \ref{lem:adjoints-of-finitely-many-normal-matrices-polynomial}, choose $p$ with $p(A)=A^*$ and $p(B)=B^*$. Thus $A^*X=XB^*$.
\item[$(\Leftarrow)$] Assume $A^*X=XB^*$. Apply the same argument to the normal matrices $A^*$ and $B^*$, choosing $q$ with $q(A^*)=A$ and $q(B^*)=B$. Hence $AX=XB$.
\end{enumerate}
\hlt{Structural Commentary}: An intertwining relation passes through every polynomial in the two normal matrices.\\
\hlt{Key Insight}: Normality makes adjoints polynomial, so intertwining transfers to adjoints.
\end{proof}
\end{theorem}

\begin{theorem} \label{thm:brauer-rank-one-eigenvalue-shift} \hlt{Brauer Rank-One Eigenvalue Shift}\\
Let $A\in\operatorname{Mat}(n,n,\mathbb{C})$, let $A\mathbf{x}=\lambda\mathbf{x}$ with $\mathbf{x}\neq\mathbf{0}$, and let $\mathbf{y}\in\mathbb{C}^n$. If $\lambda,\lambda_2,\ldots,\lambda_n$ are the eigenvalues of $A$ with multiplicity, then the eigenvalues of $A+\mathbf{x}\mathbf{y}^*$ are $\lambda+\mathbf{y}^*\mathbf{x},\lambda_2,\ldots,\lambda_n$.\\
Moreover, $\mathbf{x}$ is an eigenvector of $A+\mathbf{x}\mathbf{y}^*$ corresponding to $\lambda+\mathbf{y}^*\mathbf{x}$.
\begin{proof}
Put $\mathbf{u}=\mathbf{x}/\norm{\mathbf{x}}_2$. By Schur triangularisation, choose a unitary $U=[\mathbf{u}\ U_2]$ such that $U^*AU=\begin{bmatrix}\lambda&*\\0&T'\end{bmatrix}$, where the diagonal entries of $T'$ are $\lambda_2,\ldots,\lambda_n$. Also $U^*(\mathbf{x}\mathbf{y}^*)U
=
\begin{bmatrix}
\mathbf{y}^*\mathbf{x}&*\\
0&0
\end{bmatrix}$.
Therefore $U^*(A+\mathbf{x}\mathbf{y}^*)U$ is upper triangular with diagonal entries $\lambda+\mathbf{y}^*\mathbf{x},\lambda_2,\ldots,\lambda_n$. Finally, $(A+\mathbf{x}\mathbf{y}^*)\mathbf{x}=A\mathbf{x}+\mathbf{x}(\mathbf{y}^*\mathbf{x})=(\lambda+\mathbf{y}^*\mathbf{x})\mathbf{x}$.
\hlt{Structural Commentary}: A rank-one eigenvector update changes only the first Schur diagonal entry.\\
\hlt{Key Insight}: Align the update with the first coordinate; only one diagonal term shifts.
\end{proof}
\end{theorem}

\begin{figure}[H]
  \centering
  \begin{tikzpicture}[
  >=Stealth,
  every node/.style={font=\small},
  mat/.style={draw=gray!55, fill=gray!5, rounded corners=2pt, align=center, inner sep=6pt},
  shiftbox/.style={draw=blue!60!black, fill=blue!7, rounded corners=2pt, align=center, inner sep=6pt},
  arr/.style={->, line width=0.75pt, blue!70!black}
]
  \node[mat] (schur) at (0,0)
    {$U^*AU=\begin{bmatrix}\lambda&*\\0&T'\end{bmatrix}$};
  \node[shiftbox] (rank) at (4,0)
    {$U^*\mathbf{x}\mathbf{y}^*U=\begin{bmatrix}\mathbf{y}^*\mathbf{x}&*\\0&0\end{bmatrix}$};
  \node[mat] (sum) at (8.2,0)
    {$\begin{bmatrix}\lambda+\mathbf{y}^*\mathbf{x}&*\\0&T'\end{bmatrix}$};

  \draw[arr] (schur) -- (rank);
  \draw[arr] (rank) -- (sum);
\end{tikzpicture}

  \caption{Brauer's update shifts one Schur diagonal entry.}
  \label{fig:brauer-rank-one-shift}
\end{figure}

\begin{corollary} \label{cor:damped-rank-one-eigenvalue-reset} \hlt{Damped Rank-One Eigenvalue Reset}\\
Let $A\in\operatorname{Mat}(n,n,\mathbb{C})$ have eigenpair $(\lambda,\mathbf{x})$, let $\mathbf{y}^*\mathbf{x}=1$, and let $\tau\in\mathbb{C}$. If $\lambda,\lambda_2,\ldots,\lambda_n$ are the eigenvalues of $A$, then $A_{\tau}\defeq \tau A+(1-\tau)\lambda\mathbf{x}\mathbf{y}^*$ has eigenvalues $\lambda,\tau\lambda_2,\ldots,\tau\lambda_n$.
\begin{proof}
The eigenvalues of $\tau A$ are $\tau\lambda,\tau\lambda_2,\ldots,\tau\lambda_n$. Since $A_{\tau}=\tau A+\mathbf{x}((1-\overline{\tau})\overline{\lambda}\mathbf{y})^*$, Theorem \ref{thm:brauer-rank-one-eigenvalue-shift} shifts only $\tau\lambda$ by $(1-\tau)\lambda\mathbf{y}^*\mathbf{x}=(1-\tau)\lambda$. Thus $\tau\lambda$ becomes $\lambda$, while the other eigenvalues stay $\tau\lambda_i$.\\
\hlt{Structural Commentary}: Damping scales the spectrum; Brauer's update restores one eigenvalue.\\
\hlt{Key Insight}: Brauer's theorem changes one spectral coordinate without moving the rest.
\end{proof}
\end{corollary}

\begin{proposition} \label{prop:schur-complement-lower-right-determinant} \hlt{Schur Complement of the Lower-Right Block}\\
Let $M=
\begin{bmatrix}
A & B\\
C & D
\end{bmatrix}$, $A\in\operatorname{Mat}(r,r,\mathbb{F})$, $D\in\operatorname{Mat}(s,s,\mathbb{F})$.\\
If $D$ is invertible, then the \hlt{Schur complement of $D$ in $M$} is $M/D\defeq A-BD^{-1}C$, and $\det(M)=\det(D)\det(M/D)$.
\begin{proof}
Block Gaussian elimination gives $\begin{bmatrix}
I_r & -BD^{-1}\\
0 & I_s
\end{bmatrix}
\begin{bmatrix}
A & B\\
C & D
\end{bmatrix}
=
\begin{bmatrix}
A-BD^{-1}C & 0\\
C & D
\end{bmatrix}$.\\
The elimination matrix has determinant $1$, and the matrix on the right is lower block-triangular, so Theorem \ref{thm:block-upper-triangular-determinant} applied after transposition gives $\det(M)=\det(A-BD^{-1}C)\det(D)$.\\
\hlt{Structural Commentary}: This is the lower-right version of Schur complementation. Eliminating the upper-right block isolates the determinant contribution of $D$.\\
\hlt{Key Insight}: Eliminate by blocks, then read the determinant from the diagonal blocks.
\end{proof}
\end{proposition}

\begin{theorem} \label{thm:sylvester-determinant-identity} \hlt{Sylvester Determinant Identity}\\
Let $A\in\operatorname{Mat}(n,m,\mathbb{F})$ and $B\in\operatorname{Mat}(m,n,\mathbb{F})$. Then $\det(I_n+AB)=\det(I_m+BA)$.
\begin{proof}
Let $M=
\begin{bmatrix}
I_n & -A\\
B & I_m
\end{bmatrix}$.\\
Using Proposition \ref{prop:schur-complement-lower-right-determinant} with lower-right block $I_m$ gives $\det(M)=\det(I_n+AB)$. Using Proposition \ref{prop:schur-complement-determinant-formula} with upper-left block $I_n$ gives $\det(M)=\det(I_m+BA)$. Hence the two determinants are equal.\\
\hlt{Structural Commentary}: Two Schur complements of one block matrix give Sylvester's identity.\\
\hlt{Key Insight}: Compare the same determinant after eliminating different identity blocks.
\end{proof}
\end{theorem}

\begin{theorem} \label{thm:ab-ba-nonzero-eigenvalues} \hlt{Nonzero Eigenvalues of $AB$ and $BA$}\\
Let $A\in\operatorname{Mat}(m,n,\mathbb{C})$, $B\in\operatorname{Mat}(n,m,\mathbb{C})$, and suppose $n\geq m$. Then $p_{BA}(z)=z^{n-m}p_{AB}(z)$.\\
Consequently, for $\lambda\neq0$, $m_{AB}(\lambda)=m_{BA}(\lambda)$. If $m_{AB}(0)=k$, then $m_{BA}(0)=k+n-m$.
\begin{proof}
For $z\neq0$, Sylvester's determinant identity gives
\begin{align}
\det(zI_n-BA) =z^n\det(I_n-z^{-1}BA) =z^n\det(I_m-z^{-1}AB) =z^{n-m}\det(zI_m-AB).
\nonumber
\end{align}
Both sides are polynomials in $z$, so the identity holds for all $z$. The claimed multiplicities follow by factoring $p_{AB}$ and comparing with $p_{BA}=z^{n-m}p_{AB}$.\\
\hlt{Structural Commentary}: Sylvester's determinant identity moves the rectangular product across the determinant.\\
\hlt{Key Insight}: Passing from $AB$ to $BA$ only adds the necessary extra zero roots when the sizes differ.
\end{proof}
\end{theorem}

\begin{theorem} \label{thm:ab-ba-same-nonzero-jordan-blocks} \hlt{Nonzero Jordan Blocks of $AB$ and $BA$}\\
Let $A,B\in\operatorname{Mat}(n,n,\mathbb{C})$. For every nonzero $\lambda\in\mathbb{C}$ and every $k\geq1$, the Jordan forms of $AB$ and $BA$ contain the same number of blocks $J_k(\lambda)$.
\begin{proof}
Set
\begin{align}
X=\begin{bmatrix}AB&A\\0&0\end{bmatrix},
\qquad
Y=\begin{bmatrix}0&A\\0&BA\end{bmatrix}.
\nonumber
\end{align}
Then
\begin{align}
\begin{bmatrix}I_n&0\\B&I_n\end{bmatrix}
X
\begin{bmatrix}I_n&0\\-B&I_n\end{bmatrix}
=Y,
\nonumber
\end{align}
so $X$ and $Y$ are similar. For $\lambda\neq0$ and $p\geq1$,
\begin{align}
(X-\lambda I_{2n})^p
&=
\begin{bmatrix}
(AB-\lambda I_n)^p&*\\
0&(-\lambda)^pI_n
\end{bmatrix}, \nonumber\\
(Y-\lambda I_{2n})^p
&=
\begin{bmatrix}
(-\lambda)^pI_n&*\\
0&(BA-\lambda I_n)^p
\end{bmatrix}.
\nonumber
\end{align}
The invertible diagonal block $(-\lambda)^pI_n$ can be removed by block row and column operations, so
\begin{align}
\operatorname{rank}\bigl((X-\lambda I_{2n})^p\bigr)
&=n+\operatorname{rank}\bigl((AB-\lambda I_n)^p\bigr), \nonumber\\
\operatorname{rank}\bigl((Y-\lambda I_{2n})^p\bigr)
&=n+\operatorname{rank}\bigl((BA-\lambda I_n)^p\bigr).
\nonumber
\end{align}
Similarity of $X$ and $Y$ makes the left sides equal. Hence the rank drops defining the Weyr characteristic of $AB$ and $BA$ at $\lambda$ agree, and Theorem \ref{thm:jordan-block-sizes-from-rank-drops} gives the same nonzero Jordan blocks.\\
\hlt{Structural Commentary}: Embed $AB$ and $BA$ into similar block matrices, then use $\lambda\neq0$ to split off an invertible block.\\
\hlt{Key Insight}: Nonzero Jordan structure is detected by shifted rank drops, and those rank drops are forced to match.
\end{proof}
\end{theorem}

\begin{figure}[H]
  \centering
  \begin{tikzpicture}[
  font=\small,
  box/.style={draw=blue!60!black, fill=blue!6, rounded corners=2pt, align=center, inner sep=6pt, minimum width=2.9cm},
  aux/.style={draw=gray!65, fill=gray!6, rounded corners=2pt, align=center, inner sep=6pt, minimum width=3.2cm},
  arr/.style={-{Stealth[length=2mm]}, thick, gray!75}
]
  \node[aux] (x) at (0,0) {$X-\lambda I_{2n}$};
  \node[aux] (y) at (4.8,0) {$Y-\lambda I_{2n}$};
  \node[box] (ab) at (0,-1.55) {$AB-\lambda I_n$};
  \node[box] (ba) at (4.8,-1.55) {$BA-\lambda I_n$};

  \draw[arr, <->] (x) -- node[above] {$\sim$} (y);
  \draw[arr] (x) -- (ab);
  \draw[arr] (y) -- (ba);
\end{tikzpicture}

  \caption{For $\lambda\neq0$, the auxiliary block similarity makes the shifted ranks of $AB$ and $BA$ agree.}
  \label{fig:ab-ba-nonzero-jordan-blocks}
\end{figure}

\begin{corollary} \label{cor:ab-ba-similar-rank-power-criterion} \hlt{Similarity Criterion for $AB$ and $BA$}\\
Let $A,B\in\operatorname{Mat}(n,n,\mathbb{C})$. Then $AB$ is similar to $BA$ if and only if
\begin{align}
\operatorname{rank}\bigl((AB)^p\bigr)=\operatorname{rank}\bigl((BA)^p\bigr),\qquad p=1,\ldots,n.
\nonumber
\end{align}
\begin{proof}
If $AB$ and $BA$ are similar, then their powers are similar and have equal rank.\\
Conversely, Theorem \ref{thm:ab-ba-same-nonzero-jordan-blocks} gives equality of all nonzero Jordan blocks. The displayed rank equalities, together with the trivial equality at $p=0$, give the same Weyr characteristic at $\lambda=0$ by Theorem \ref{thm:jordan-block-sizes-from-rank-drops}. Hence the nilpotent Jordan blocks also agree, so the two Jordan forms are identical up to block order.\\
\hlt{Structural Commentary}: Nonzero blocks agree automatically; only the nilpotent blocks must be tested.\\
\hlt{Key Insight}: Powers of $AB$ and $BA$ measure exactly the zero-Jordan-block data.
\end{proof}
\end{corollary}

\begin{method} \label{meth:checking-ab-ba-similarity} \hlt{Checking Whether $AB$ and $BA$ Are Similar}\\
For $A,B\in\operatorname{Mat}(n,n,\mathbb{C})$:
\begin{enumerate}[label=(\roman*)]
\setlength{\itemsep}{0pt}
\item The nonzero Jordan blocks of $AB$ and $BA$ already match.
\item Compute $\operatorname{rank}\bigl((AB)^p\bigr)$ and $\operatorname{rank}\bigl((BA)^p\bigr)$ for $p=1,\ldots,n$.
\item If all ranks agree, then $AB\sim BA$; if some rank differs, then only the nilpotent Jordan structure differs.
\end{enumerate}
\end{method}

\begin{corollary} \label{cor:hermitian-products-and-conjugate-products} \hlt{Hermitian Products and Conjugate Products}\\
Let $A,H,K\in\operatorname{Mat}(n,n,\mathbb{C})$, and suppose $H$ and $K$ are Hermitian. Then $HK$ is similar to $KH$, and $A\overline{A}$ is similar to $\overline{A}A$.
\begin{proof}
For each $p=1,\ldots,n$,
\begin{align}
\operatorname{rank}\bigl((HK)^p\bigr)
=\operatorname{rank}\bigl(((HK)^p)^*\bigr)
=\operatorname{rank}\bigl((KH)^p\bigr).
\nonumber
\end{align}
Also,
\begin{align}
\operatorname{rank}\bigl((A\overline{A})^p\bigr)
=\operatorname{rank}\bigl(\overline{(A\overline{A})^p}\bigr)
=\operatorname{rank}\bigl((\overline{A}A)^p\bigr).
\nonumber
\end{align}
Apply Corollary \ref{cor:ab-ba-similar-rank-power-criterion} in both cases.\\
\hlt{Structural Commentary}: Adjoint and conjugation preserve rank while reversing the product order in these two situations.\\
\hlt{Key Insight}: The rank-power test turns order reversal into similarity.
\end{proof}
\end{corollary}

\begin{theorem} \label{thm:conjugate-similarity-hermitian-product-equivalences} \hlt{Conjugate Similarity and Hermitian Products}\\
Let $A\in\operatorname{Mat}(n,n,\mathbb{C})$. The following are equivalent:
\begin{enumerate}[label=(\roman*)]
\setlength{\itemsep}{0pt}
\item $A$ is similar to $\overline{A}$.
\item $A$ is similar to $A^*$.
\item $A$ is similar to $A^*$ via a Hermitian similarity matrix.
\item $A=HK=LM$, where $H,K,L,M$ are Hermitian, and $H,M$ are invertible.
\item $A=HK$ for some Hermitian matrices $H,K$.
\end{enumerate}
\begin{proof}
\begin{enumerate}[label=\roman*.]
\setlength{\itemsep}{0pt}
\item[$(i)\Rightarrow(ii)$] Theorem \ref{thm:matrix-similar-to-transpose-and-symmetric-factorizations} applied to $\overline{A}$ gives $\overline{A}\sim(\overline{A})^t=A^*$. Thus $A\sim A^*$.
\item[$(ii)\Rightarrow(iii)$] Choose invertible $S$ such that $A=SA^*S^{-1}$, so $AS=SA^*$. For $\theta\in\mathbb{R}$, put $S_{\theta}=e^{i\theta}S$. Then $AS_{\theta}=S_{\theta}A^*$ and, by taking adjoints, $AS_{\theta}^*=S_{\theta}^*A^*$. Hence
\begin{align}
A(S_{\theta}+S_{\theta}^*)=(S_{\theta}+S_{\theta}^*)A^*.
\nonumber
\end{align}
The matrix $H_{\theta}=S_{\theta}+S_{\theta}^*$ is Hermitian. Also
\begin{align}
S_{\theta}^{-1}H_{\theta}=I_n+e^{-2i\theta}S^{-1}S^*.
\nonumber
\end{align}
Only finitely many values of $-e^{2i\theta}$ can be eigenvalues of $S^{-1}S^*$, so choose $\theta$ for which $H_{\theta}$ is invertible. Then $A=H_{\theta}A^*H_{\theta}^{-1}$.
\item[$(iii)\Rightarrow(iv)$] Let $R$ be invertible Hermitian and satisfy $A=RA^*R^{-1}$. Then $AR=RA^*$ and $R^{-1}A=A^*R^{-1}$. Thus $A^*R^{-1}$ and $RA^*$ are Hermitian, and
\begin{align}
A=R(A^*R^{-1})=(RA^*)R^{-1}.
\nonumber
\end{align}
This gives the two asserted Hermitian factorizations, with the first factor in the first product and the second factor in the second product invertible.
\item[$(iv)\Rightarrow(v)$] Immediate.
\item[$(v)\Rightarrow(i)$] If $A=HK$ with $H,K$ Hermitian, then Corollary \ref{cor:hermitian-products-and-conjugate-products} gives $HK\sim KH$. By Theorem \ref{thm:matrix-similar-to-transpose-and-symmetric-factorizations},
\begin{align}
KH\sim (KH)^t=H^tK^t=\overline{H}\,\overline{K}=\overline{A}.
\nonumber
\end{align}
Hence $A\sim\overline{A}$.
\end{enumerate}
\hlt{Structural Commentary}: The equivalences move between conjugation, adjoints, and Hermitian products by combining transpose similarity with a Hermitianized similarity matrix.\\
\hlt{Key Insight}: A phase-rotated average $e^{i\theta}S+e^{-i\theta}S^*$ turns an arbitrary similarity to $A^*$ into a Hermitian one.
\end{proof}
\end{theorem}

\begin{figure}[H]
  \centering
  \begin{tikzpicture}[
  font=\small,
  node/.style={draw=blue!60!black, fill=blue!6, rounded corners=2pt, align=center, inner sep=6pt, minimum width=2.9cm},
  mid/.style={draw=red!60!black, fill=red!6, rounded corners=2pt, align=center, inner sep=6pt, minimum width=3.3cm},
  arr/.style={-{Stealth[length=2mm]}, thick, gray!75},
  lbl/.style={fill=white, text=gray!70!black, inner xsep=2pt, inner ysep=1pt}
]
  \node[node] (conj) at (0,0) {$A\sim\overline{A}$};
  \node[node] (adj) at (5.0,0) {$A\sim A^*$};
  \node[mid] (herm) at (10.4,0) {$A=HA^*H^{-1}$\\$H=H^*$};
  \node[node] (prod) at (5.0,-2.2) {$A=HK$\\$H=H^*,\ K=K^*$};

  \draw[arr] (conj.east) -- node[lbl, above, yshift=2pt] {transpose} (adj.west);
  \draw[arr] (adj.east) -- node[lbl, above, yshift=2pt] {phase average} (herm.west);
  \draw[arr] (herm.south west) -- node[lbl, right, xshift=2pt] {factor} (prod.east);
  \draw[arr] (prod.west) -- node[lbl, left, xshift=-2pt] {$HK\sim KH$} (conj.south east);
\end{tikzpicture}

  \caption{Conjugate similarity, adjoint similarity, and Hermitian product factorisations are equivalent.}
  \label{fig:conjugate-similarity-hermitian-products}
\end{figure}

\begin{proposition} \label{prop:block-centrosymmetric-commutation} \hlt{Block Centrosymmetry and Commutation}\\
Let $M=
\begin{bmatrix}
A & B\\
C & D
\end{bmatrix}\in\operatorname{Mat}(2n,2n,\mathbb{F})$,
$L_{2n}=
\begin{bmatrix}
0 & I_n\\
I_n & 0
\end{bmatrix}$.\\
Then $M$ has block centrosymmetric form $A=D$ and $B=C$ if and only if $L_{2n}M=ML_{2n}$.
\begin{proof}
Block multiplication gives $L_{2n}M=
\begin{bmatrix}
C & D\\
A & B
\end{bmatrix}$ and 
$ML_{2n}=
\begin{bmatrix}
B & A\\
D & C
\end{bmatrix}$.\\
The two matrices are equal exactly when $A=D$ and $B=C$.\\
\hlt{Structural Commentary}: The matrix $L_{2n}$ swaps the two block coordinates. Commuting with that swap means $M$ treats the two block positions symmetrically.\\
\hlt{Key Insight}: Compute both block products; equality forces matching opposite blocks.
\end{proof}
\end{proposition}

\begin{theorem} \label{thm:frobenius-rank-inequality} \hlt{Frobenius Rank Inequality}\\
Let $A\in\operatorname{Mat}(m,k,\mathbb{F})$, $B\in\operatorname{Mat}(k,p,\mathbb{F})$, and $C\in\operatorname{Mat}(p,n,\mathbb{F})$. $\operatorname{rank}(AB)+\operatorname{rank}(BC) \leq \operatorname{rank}(B)+\operatorname{rank}(ABC)$.
\begin{proof}
Let $R=\operatorname{col}(B)$ and $S=\operatorname{col}(BC)$. Then $S\subseteq R$. The induced map $T_A$ maps $R$ onto $\operatorname{col}(AB)$ and maps $S$ onto $\operatorname{col}(ABC)$. Rank-nullity on the restrictions $T_A|_R$ and $T_A|_S$ gives
\begin{align}
\operatorname{rank}(AB)&=\dim R-\dim(R\cap\ker A), \nonumber\\
\operatorname{rank}(ABC)&=\dim S-\dim(S\cap\ker A). \nonumber
\end{align}
Since $S\cap\ker A\subseteq R\cap\ker A$, we have $\dim(S\cap\ker A)\leq\dim(R\cap\ker A)$.\\
Substituting $\dim R=\operatorname{rank}(B)$ and $\dim S=\operatorname{rank}(BC)$ gives the stated inequality.\\
\hlt{Structural Commentary}: Frobenius tracks $A$ on nested column spaces.\\
\hlt{Key Insight}: The smaller column space cannot lose more dimensions to $\ker A$ than the larger one.
\end{proof}
\end{theorem}

\begin{corollary} \label{cor:frobenius-rank-inequality-special-cases} \hlt{Rank Inequalities from Frobenius}\\
The following are special cases of Theorem \ref{thm:frobenius-rank-inequality}:
\begin{enumerate}[label=(\roman*)]
\setlength{\itemsep}{0pt}
\item
$\operatorname{rank}(AB)\leq\min(\operatorname{rank}(A),\operatorname{rank}(B))$;
\item
$\operatorname{rank}(A)+\operatorname{rank}(B)-k\leq\operatorname{rank}(AB)$.
\end{enumerate}
\begin{proof}
\begin{enumerate}[label=(\roman*)]
\setlength{\itemsep}{0pt}
\item
Apply Frobenius to generic factors $X,Y,Z$. Taking $Z=0$ gives $\operatorname{rank}(XY)\leq\operatorname{rank}(Y)$, and taking $X=0$ gives $\operatorname{rank}(YZ)\leq\operatorname{rank}(Y)$. Renaming the factors gives $\operatorname{rank}(AB)\leq\operatorname{rank}(B)$ and $\operatorname{rank}(AB)\leq\operatorname{rank}(A)$.
\item
Taking the middle matrix to be $I_k$ gives $\operatorname{rank}(A)+\operatorname{rank}(B)\leq k+\operatorname{rank}(AB)$, which is Sylvester's inequality.
\end{enumerate}
\hlt{Structural Commentary}: Frobenius is the three-factor parent inequality. The usual upper and lower product-rank bounds appear by making one factor zero or identity.\\
\hlt{Key Insight}: Zero factors isolate upper bounds; an identity middle factor yields the Sylvester lower bound.
\end{proof}
\end{corollary}

\begin{proposition} \label{prop:sylvester-law-of-nullity} \hlt{Sylvester's Law of Nullity}\\
For a matrix $X$, let $\operatorname{nullity}(X)=\dim(\ker X)$; this is the dimension of its nullspace.
\begin{enumerate}[label=(\roman*)]
\setlength{\itemsep}{0pt}
\item If $A\in\operatorname{Mat}(m,k,\mathbb{F})$ and $B\in\operatorname{Mat}(k,n,\mathbb{F})$, then $\operatorname{nullity}(AB)\leq\operatorname{nullity}(A)+\operatorname{nullity}(B)$
\item
If $A,B\in\operatorname{Mat}(n,n,\mathbb{F})$, then $\max(\operatorname{nullity}(A),\operatorname{nullity}(B))\leq\operatorname{nullity}(AB)$
\end{enumerate}
\begin{proof}
\begin{enumerate}[label=(\roman*)]
\setlength{\itemsep}{0pt}
\item
By rank-nullity and Theorem \ref{thm:sylvester-rank-inequality},
\begin{align*}
\operatorname{nullity}(AB) &=n-\operatorname{rank}(AB) \leq n-\operatorname{rank}(A)-\operatorname{rank}(B)+k \\
&=(k-\operatorname{rank}(A))+(n-\operatorname{rank}(B)) =\operatorname{nullity}(A)+\operatorname{nullity}(B).
\end{align*}
\item For square $A,B$, $\operatorname{rank}(AB)\leq\operatorname{rank}(A)$, $\operatorname{rank}(AB)\leq\operatorname{rank}(B)$. Subtract from $n$ gives desired lower bounds.
\end{enumerate}
\hlt{Structural Commentary}: Nullity is rank loss from the domain.\\
\hlt{Key Insight}: Use rank-nullity to translate rank bounds into kernel bounds.
\end{proof}
\end{proposition}
